\documentclass[11pt]{article}
\usepackage[margin=25mm]{geometry}
\usepackage{fontspec}
\setmainfont{Latin Modern Roman}
\usepackage{amsmath,amssymb,mathtools}
\usepackage{unicode-math}
\setmathfont{Latin Modern Math}
\usepackage{xurl,hyperref}
\hypersetup{colorlinks=true,urlcolor=blue}
\setcounter{secnumdepth}{0}
\setlength{\parindent}{0pt}
\setlength{\parskip}{5pt}
\setlength{\emergencystretch}{3em}
\newcommand{\sref}[2]{\hyperlink{source:#1:#2}{[S#2]}}
\newcommand{\cataloguescope}[1]{\paragraph{Recorded target.}#1}
\newcommand{\R}{\mathbb R}
\newcommand{\C}{\mathbb C}
\newcommand{\N}{\mathbb N}
\newcommand{\Z}{\mathbb Z}
\newcommand{\Q}{\mathbb Q}
\newcommand{\F}{\mathbb F}
\newcommand{\eps}{\varepsilon}
\begin{document}
\section{Infrared Equivalence of Stack Presentations}\label{statement}

\subsection{Definitions and mathematical statement}

\paragraph{Shared notation.}$\mathbb N=\{1,2,\ldots\}$; the standard number fields have their usual meanings. All additional conventions are stated locally.


\paragraph{Definitions and conventions.} For a physically realizable presentation $[X/G]$ of a stack, the associated worldsheet theory is the report's $G$-gauged sigma model on $X$. Renormalization-group flow changes its couplings with scale; its infrared theory describes the long-distance limit. Two infrared theories being the same means physical equivalence in the report's sense, not equality of their microscopic fields, gauge groups, or bare actions.

\paragraph{Source scope and qualifications.} The source restricts to physically realizable presentations and formulates infrared equivalence as a conjecture. It neither supplies a fully defined universal RG flow nor asserts equivalence for arbitrary stack presentations; the qualification is essential.

\paragraph{Catalogue mathematical target.}
\paragraph{Catalogue wording retained for traceability.}
Physical theories associated with different physically realizable presentations of the same stack should become the same infrared theory under renormalization-group flow.

\paragraph{Saved literature context.} The supplied catalogue's assessment dated 2026-08-21 records: Presentation independence under worldsheet renormalization-group flow is supported by physical consistency checks and examples, but the conjectured equality of the infrared theories is not established in general. The assessment is retained as dated provenance; it is not a new comprehensive literature or open-status certification.
\subsection{Short English statement}
Show that different physically realizable quotient presentations of one stack lead to equivalent long-distance string theories.
\subsection{Research attempt}
\paragraph{Model, date and investigated scope.}
OpenAI GPT-6 Astra (configured model \texttt{gpt-6-astra}), ultra reasoning; generated 11 October 2026. Establishes presentation invariance of a finite groupoid partition sum and diagnoses the missing gauge-volume normalization.

\paragraph{Sources and literature check.}
The complete supplied definition, including source qualifications and the dated literature assessment, was reviewed. The calculation below is an elementary restricted attempt motivated by [S2]; it does not claim novelty or a fresh literature-status determination.

\paragraph{Context and investigated scope.} The complete supplied definition, including source qualifications and the dated literature assessment, was reviewed. The calculation below is an elementary restricted attempt motivated by [S2]; it does not claim novelty or a fresh literature-status determination.

\paragraph{An exact finite presentation check.} A finite groupoid has finitely many objects and morphisms. Its weighted cardinality is $|\mathcal G|=\sum_{[x]}1/|\operatorname{Aut}(x)|$, the sum over isomorphism classes. If a finite group $G$ acts on a finite set $X$, the action groupoid has
\[
 |X//G|=\sum_{\text{orbits }O}1/|G_x|
 =\sum_O|O|/|G|=|X|/|G|.
\]
The middle equality follows from the bijection $G/G_x\to O$. An equivalence of finite groupoids bijects isomorphism classes and gives isomorphic automorphism groups, so this weighted cardinality is presentation independent.

\paragraph{A concrete change of presentation.} Let a finite group $H$ act freely on itself and trivially on $X$. Then $[X\times H/(G\times H)]$ presents the same finite quotient as $[X/G]$, with the obvious independent actions. The unnormalized counts differ by $|H|$, but the normalized zero-dimensional partition sum agrees:
\[
 \frac{|X\times H|}{|G\times H|}=\frac{|X|}{|G|}.
\]
The same proof works for any weight constant on groupoid isomorphism classes and pulled back under the equivalence.

\paragraph{Why this is not an RG theorem.} This exact calculation checks the gauge-volume normalization in a finite toy model. A worldsheet sigma model has fields, kinetic terms, possible anomalies and scale-dependent couplings, none of which is represented by the count. Stack equivalence alone does not identify the chosen measures and renormalized actions. The first unresolved step is a correspondence between physically realizable presentations preserving those data along a defined RG flow. No conclusion for arbitrary presentations, and no infrared equivalence for the source theories, follows from the counting identity.

\paragraph{Verification and adversarial audit.}
Checked orbit-stabilizer normalization, free extra-group presentation and invariant weights; restricted the result to zero-dimensional finite models. Model self-audit only; no independent human proof checking. Prompts A and B were applied to this canonical record and its qualified source scope. Checked orbit-stabilizer normalization, free extra-group presentation and invariant weights; restricted the result to zero-dimensional finite models. Retained unresolved status and zero original-target progress; No equivalence of worldsheet measures, renormalized couplings and infrared limits is constructed.

Coordinator checked orbit-stabilizer weighted cardinality and the free extra-group presentation, without inferring an RG equivalence. This additional check was model-assisted, not independent human verification.

\paragraph{Tools, process and contributions.}
Python 3 was used for file preparation and structural checks; no numerical experiment or theorem prover was used. Contribution prompt A was followed by prompt B. GPT-6 Astra with ultra reasoning drafted the research and performed its stated self-audit; any additional coordinator contribution is identified in the verification and audit. A GPT-6-based Codex coordinator assembled the manuscript and checked repository integration; its exact variant and reasoning configuration were not exposed. No human mathematical contribution, independent human proof check or formal proof certification is supplied. Complete derivations and the consulted source references are included here; private reasoning transcripts are not published.

\paragraph{Final status and first remaining gap.}
\textbf{UNRESOLVED}. No equivalence of worldsheet measures, renormalized couplings and infrared limits is constructed.

\section{Completion Estimate}
\noindent\textbf{COMPLETION: 0\%}
The calculation is an auxiliary or restricted attempt and establishes no newly verified portion of the full source target. No equivalence of worldsheet measures, renormalized couplings and infrared limits is constructed.

\subsection{Sources}
\begin{itemize}
\item[\hypertarget{source:statement:1}{S1}] ULAM.AI and the UnsolvedMath Contributors, UnsolvedMath, catalogue entry. Catalogue curation under CC BY 4.0; definitions and English exposition added locally.
\url{https://huggingface.co/datasets/ulamai/UnsolvedMath}.
\item[\hypertarget{source:statement:2}{S2}] Mini-Workshop: Heterotic Strings, Derived Categories, and Stacks (2005), Oberwolfach Reports, DOI 10.4171/owr/2005/53, pp. 3019--3060; official report source and mathematical context.
\url{https://doi.org/10.4171/owr/2005/53}.
\item[\hypertarget{source:statement:3}{S3}] Official EMS Press PDF of the same Oberwolfach report; PDF page 19, printed page 3037.
\url{https://ems.press/content/serial-article-files/46027}.
\item[\hypertarget{source:statement:4}{S4}] stack GLSM analysis; reference in the saved dated catalogue literature review.
\url{https://arxiv.org/abs/hep-th/0502053}.
\end{itemize}
\label{end:statement}
\end{document}
